Nguồn: \href{https://www.facebook.com/notes/642263246455963/}{FB tác giả}

Hôm nay, người thân, học trò, những người quý mến GS Phan Đình Diệu đã đến viếng và đưa ông trở về với cát bụi. Ở xa, không về được, tôi viết bài này như một lời tiễn đưa giáo sư,  một trong rất ít người mà tôi thực sự ngưỡng mộ. Cả về trí tuệ lẫn nhân cách!

 

GS Phan Đình Diệu sinh năm 1936, mất ngày 13 tháng 5 năm 2018. Tên tuổi và tài năng của ông, xã hội biết; nhân cách của ông, người đời kính trọng. Đọc những gì mà mấy hôm nay mọi người viết về ông thì sẽ rõ. Người vừa có trí, vừa có nhân như GS Phan Đình Diệu, giờ có không nhiều.

 

Thiên hạ viết về GS Phan Đình Diệu khá nhiều. Ở đây, tôi muốn viết về ông dưới con mắt của mình, qua những gì mình được chứng kiến. Viết về con người ông như tôi vẫn biết.

 

Tôi có may mắn, ngày còn ở nhà, nhiều lần được gặp giáo sư, được chứng kiến những câu chuyện thân tình, những trăn trở của ông  về đất nước cũng như nhiều chuyện khác. Bố tôi và GS Phan Đình Diệu là đôi bạn có thể nói là tâm giao, tri kỉ. Một người ngành văn, một người ngành toán, hai lĩnh vực hoàn toàn khác nhau nhưng có cái gì đó trong họ tương đồng, gắn kết lại với nhau. Họ dành cho nhau những tình cảm thật đặc biệt. 

 

Ngày trước lúc cùng khu tập thể, hay sau đó khi gia đình giáo sư chuyển ra khu Đội Cấn bố hay nhờ đèo qua thăm ông. Hai người hay qua lại, đàm đạo với nhau, và vì thế tôi cũng được “hóng hớt” nghe chuyện hai người.

 

GS Phan Đình Diệu là người thông minh hiếm có. Ông tuộc tuýp người đa tài, quan tâm nhiều lĩnh vực. Ở ông, sự thông minh lộ toát ra ngoài, gây ấn tượng cho nhiều người chỉ qua một lần gặp gỡ. Tôi không biết người Việt có ai ngoài ông, chỉ trong vòng 5 năm (kể cả thời gian học tiếng) mà bảo vệ thành công cả PTS rồi TS (tương đương TSKH của Việt Nam) của Liên Xô? Nhiều giáo sư Nga để bảo vệ TS mất cả hàng chục năm. Ông không chỉ giỏi toán mà còn làm thơ khá hay. Ông quan tâm đến triết học, khoa học xã hội và có nhiều phát biểu khiến giới KHXH phải nể. Tôi còn nhớ dạo những năm 80, khi cậu tôi GS Nguyễn Đức Đàn về làm TBT Tạp chí Nghiên cứu nghệ thuật đã cho ra những số với những bài rất hay mà các tác giả là dân KHTN. Những bài viết của GS Tạ Quang Bửu, bài của GS vật lí Cao Chi về đối xứng và phản đối xứng trong kiến trúc lăng tẩm Huế… GS Phan Đình Diệu, tôi vẫn nhớ, lúc đó cũng có hai bài thơ: một về Tagor, một về tâm sự khi 2 người con dân Việt tình cờ gặp nhau trên một sân ga nhỏ ở Đức nhưng hờ hững như người dưng. Cuối những năm 80, viện HLKH Nga có mời phía Việt Nam tham gia viết chung một tuyển tập các bài viết về triết học. Kết quả phía Việt Nam chỉ riêng bài của GS Phan Đình Diệu được chọn, các bài khác bị từ chối vì lí do .. không phải thuộc triết mà là chính trị học.

 

Tư chất thông minh của giáo sư, nhiều người đều biết. Nhưng cái tôi muốn nói, cái mình thực sự ấn tượng và nể trọng ở ông lại là nhân cách, sự tìm tòi trăn trở cho đất nước và bản lĩnh người. Phát biểu công khai, đàng hoàng, có lí có tình, đánh bài ngửa. Nhiều người coi ông là người có bất đồng chính kiến nhưng ông phê phán với ý thức xây dựng, dám làm dám chịu. Chính thái độ này, cái cách làm chính trực, cùng sự thông minh, đã chinh phục được không ít vị lãnh đạo cao cấp của đảng và nhà nước và họ đã trở thành những người “bảo trợ” cho ông, tránh được những rủi ro bất trắc. Có lẽ, ngoài Phan Đình Diệu, không mấy nhà đối lập chính kiến lại được mời làm khách đến nói chuyện tại nhà riêng của không ít vị lãnh đạo cao cấp của nhà nước Việt Nam.

 

Lần đầu được tiếp xúc với ông là khoảng đầu những năm 80 khi tôi vừa tốt nghiệp mới quay về nước. Lần đó ông kể: sau chuyến đi Mĩ về, có quan điểm bảo thủ muốn coi những người rời bỏ đất nước là kẻ phản bội đất nước, không được phép về thăm lại Việt Nam. Ông đã phải đấu tranh, thuyết phục để “cái điều ngớ ngẩn” này không được chấp nhận. Đơn thuần chỉ với lí do: đất nước là của chung của mọi người con đất Việt, không ai có thể tước đi cái quyền đó. Trước đấy không lâu, ông đã có một phát biểu động trời tại Quốc hội, công khai đòi TBT Lê Duẩn từ chức. Ông kể lại: sau hôm đó, đại tướng Võ Nguyên Giáp (người vỗn rất quí mến ông) có nhắn tin mời ông đến nhà. Khác với những lần trước, Võ Nguyên Giáp quân phục chỉnh tề, đón ông với mặt không vui, không thân tình như những mọi lần. Nhưng sau khi nghe ông trình bày lại, lí giải quan điểm của mình, đại tướng thôi căng thẳng, vui vẻ trở lại.

 

Cương trực, thẳng thắn, công khai phát biểu chính kiến của mình một cách đàng hoàng mà nhiều điều bất đồng không dễ gì chính quyền nghe được. Đó là điều không phải ai cũng dám và làm được. GS Phan Đình Diệu là một người như thế! Tôi vẫn nhớ, có lần cô Hương phải nhờ bố tôi với tư cách bạn vong niên để khuyên giáo sư nói nhẹ bớt chứ không thì nguy hiểm, nhiều cái không hay có thể đến. Nói thì nói vậy, chứ người như GS Phan Đình Diệu không dễ gì làm thay đổi chính kiến nếu không dùng lí lẽ thuyết phục được ông.

 

Những suy tư, những quan điểm đó của ông được thể hiện qua những bài viết, bài được công bố bài thì không. Mỗi lần như vậy, ông lại đưa cho bố tôi đọc, để chia sẻ và đồng cảm. Và nhờ vậy, tôi cũng được đọc theo. Thật tình, tôi rất thích đọc những bài viết của GS Phan Đình Diệu. Vừa sắc sảo trí tuệ, có tầm nhìn xa, vừa trăn trở thiết tha của một nhà trí thức đau mình về đất nước. Những bài viết này thật sự là một tài sản có giá trị, mà bố tôi vẫn nhắc các con ông nên sưu tập lại in thành sách, đừng bỏ phí.

 

Có thể kể lại rất nhiều, tuy nhiên ở đây xin chỉ nhắc lại đôi điều.

 

Hẳn mọi người còn nhớ vụ kỉ luật Trần Xuân Bách, cựu UV BCT. Bài phát biểu “Chủ nghĩa xã hội thật sự là gì?” của Trần Xuân Bách có nhiều trùng lặp với quan điểm của Phan Đình Diệu thể hiện trong bài viết của ông mà trước đó tôi đã được đọc. Sự tự diễn biến của Trần Xuân Bách có chịu nhiều ảnh hưởng của Phan Đình Diệu.

 

Tôi vẫn nhớ lần GS Tương Lai đứng ra làm trng gian mời ba người đến nói chuyện trực tiếp cho TT Phạm Văn Đồng về về xã hội, đât nước. Đó là GS Phan Đình Diệu với tư cách một nhà khoa học tự nhiên, Ma Văn Kháng với góc nhìn của một nhà văn và bố tôi - một nhà nghiên cứu văn học. Mọi người thống nhất với nhau sẽ nói thẳng, nói thật, đúng với suy nghĩ của mình. Tôi không biết sau đó thế nào nhưng được biết trước khi gặp TT Phạm Văn Đồng GS Phan Đình Diệu có buổi thuyết trình cho UB KHXH. Buổi nói chuyện gây ấn tượng khá lớn cho họ, cả về cách suy nghĩ lần bản lĩnh dám nói.  

 

Rồi cái lần nhập Bộ Giáo dục và Bộ Đại học, có ý kiến ở trên muốn đề cử GS Phan Đình Diệu làm bộ trưởng giáo dục, hay những lần có tâm sự này khác chuyện cơ quan ông cũng hay tâm sự với bố tôi.

 

Tôi nhớ GS Phan Đình Diệu có kể lại: năm 90 khi nhà văn Bùi Tín qua Pháp và muốn ở lại với ý muốn để góp ý cho ĐCS và nhà nước Việt Nam. Trước khi quyết định Bùi Tín có đến gặp và xin ý kiến GS Phan Đình Diệu. Ông có khuyên Bùi Tín đại khái như thế này: nếu anh thực sự yêu nước, muốn đóng góp thì nên về và phát biểu tại Việt Nam chứ anh ở lại và nói ở nước ngoài thì chẳng ai sẽ hiểu cho anh đâu. Không hiểu Bùi Tín nghĩ gì nhưng sau đó quyết định xin tỵ nạn ở lại Pháp.

 

Ngày trước, cùng với ĐCS lãnh đạo ở Việt Nam còn có 2 đảng khá lâu đời khác là đảng Xã hội và đảng Dân Chủ của giới trí thức. Năm 1988 do có cảm giác về vị thế thực của hai đảng này mà các vị lãnh đạo chủ chốt xin giải tán đảng. Tưởng “làm mình làm mẩy” cho vui không ngờ bị giải tán thật. Hơn nữa, các tờ báo cơ quan ngôn luận của hai đảng vì thế cũng bị đóng cửa nốt, không còn sân chơi, nơi có thể giao lưu và bày tỏ ý kiến của giới trí thức thuộc hai đảng này. Trong tâm trạng buồn, GS Nguyễn Xiển (Tổng thư kí đảng Xã hội Việt Nam) có gặp GS Phan Đình Diệu và nói, đại ý: thế hệ chúng tôi có nhiều sai lầm. Cụ Xiển vừa nói vừa ứa nước mắt. GS Phan Đình Diệu có an ủi và nói với GS Nguyễn Xiển: nhưng dù sao thế hệ các bác cũng rất đẹp. Ngày trước các bác trong sáng, mục đích rất tốt đẹp. Còn thế hệ chúng tôi thì không những bị lừa mà còn lừa dối lẫn nhau.

 

Kể lại đôi điều mà tôi được biết về GS Phan Đình Diệu là muốn phác họa lại nhân cách, cái hình ảnh của ông, như tôi vẫn biết. Một con người thông minh, tài hoa, vừa có trí vừa có nhân. “Một người kẻ sĩ chân chính và chính trực” (BBC), hết lòng vì đất nước, tử tế với bạn bè. Có lẽ vì thế các con ông đã được hưởng nhiều điều tốt đẹp từ ông.

 

Vĩnh biệt chú, GS Phan Đình Diệu, người bạn lớn của gia đình, người mà tôi hết lòng ngưỡng mộ!

 

Moscow, 18-05-2018